%-------------------------
% Academic / Industry CV — LaTeX Template
% A clean, professional CV for ML / Computer Vision engineers
% Compiler: pdfLaTeX (Overleaf-ready)
% License: MIT — free to use and modify
%-------------------------

\documentclass[11pt,a4paper]{article}

%---------- PACKAGES ----------
\usepackage[left=1.8cm, right=1.8cm, top=1.6cm, bottom=1.6cm]{geometry}
\usepackage[utf8]{inputenc}
\usepackage[T1]{fontenc}
\usepackage{lmodern}
\usepackage{enumitem}
\usepackage{titlesec}
\usepackage{xcolor}
\usepackage{hyperref}
\usepackage{fontawesome5}
\usepackage{tabularx}
\usepackage{ragged2e}

%---------- COLORS ----------
% Tip: change these to personalize your CV instantly
\definecolor{accentcolor}{HTML}{8B2500}
\definecolor{darktext}{HTML}{1a1a2e}
\definecolor{subtext}{HTML}{555555}
\definecolor{rulecolor}{HTML}{c45d3e}
\definecolor{linkblue}{HTML}{1a5276}

%---------- HYPERREF ----------
\hypersetup{
    colorlinks=true,
    linkcolor=linkblue,
    urlcolor=linkblue,
    pdfauthor={Your Name},
    pdftitle={CV - Your Name},
}

%---------- SPACING ----------
\setlength{\parindent}{0pt}
\setlength{\parskip}{0pt}

%---------- SECTION FORMATTING ----------
\titleformat{\section}{%
    \color{accentcolor}\large\bfseries\scshape
}{}{0em}{}[\vspace{2pt}{\color{rulecolor}\hrule height 0.6pt}]
\titlespacing*{\section}{0pt}{14pt}{7pt}

%---------- CUSTOM COMMANDS ----------
\newcommand{\cventry}[3]{%
    \textbf{#1}\hfill\textit{#2}\par\vspace{2pt}%
    #3\par\vspace{7pt}%
}

\newcommand{\cvproject}[4]{%
    \textbf{#1}\hfill\textit{#2}\par\vspace{2pt}%
    #3\par\vspace{1pt}%
    {\small\color{subtext}\textit{Tech: #4}}\par\vspace{7pt}%
}

\newcommand{\skillrow}[2]{%
    \textbf{#1:} & #2 \\[3pt]
}

%========================================
\begin{document}
\pagestyle{empty}
\color{darktext}

%========== HEADER ==========
\begin{center}
    {\LARGE\bfseries Alex Chen}\\[7pt]
    {\large\color{subtext} Object Detection\enspace$\cdot$\enspace Computer Vision Engineer}\\[9pt]
    {\small
        \faEnvelope\enspace \href{mailto:alex.chen.cv@outlook.com}{alex.chen.cv@outlook.com}
        \enspace$|$\enspace
        \faPhone\enspace +1\,(555)\,012\,3456
    }\\[3pt]
    {\small
        \faGithub\enspace \href{https://github.com/alexchen-cv}{github.com/alexchen-cv}
        \enspace$|$\enspace
        \faGlobe\enspace \href{https://alexchen-cv.github.io/}{alexchen-cv.github.io}
        \enspace$|$\enspace
        \faLinkedin\enspace \href{https://linkedin.com/in/alexchen-cv}{LinkedIn}
    }\\[3pt]
    {\small\color{subtext}
        Open to relocation: Germany\enspace$\cdot$\enspace Netherlands\enspace$\cdot$\enspace EU
        \enspace$|$\enspace
        Eligible for EU Blue Card
    }
\end{center}
\vspace{2pt}

%========== PROFESSIONAL SUMMARY ==========
\section{Professional Summary}

Computer vision engineer with an M.Eng.\ in Computer Science and 2+ years of research focused on real-time object detection and multi-scale feature fusion. Designed and deployed two detection architectures --- an attention-enhanced YOLO variant and a transformer-CNN hybrid detector --- achieving state-of-the-art mAP on COCO and domain-specific industrial datasets. Experienced in edge deployment (TensorRT / ONNX) and building end-to-end perception pipelines from data annotation to production inference.

%========== EDUCATION ==========
\section{Education}

\cventry{M.Eng.\ in Computer Science}{Sep 2022 -- Jun 2025}{%
    University of Science and Technology, China\par\vspace{1pt}
    Thesis: \textit{Multi-Scale Attention Networks for Real-Time Object Detection}\enspace$|$\enspace Advisor: Prof.\ Wei Zhang
    \begin{itemize}[leftmargin=14pt, nosep, topsep=3pt, itemsep=1pt]
        \item Proposed two novel detection architectures; benchmarked on COCO, VOC, and custom industrial datasets.
        \item Trained on NVIDIA A100 cluster with PyTorch; implemented full ablation studies and latency profiling.
    \end{itemize}
}

\cventry{B.Eng.\ in Software Engineering}{Sep 2018 -- Jun 2022}{%
    National University, China\par\vspace{1pt}
    Thesis: \textit{Pedestrian Detection in Crowded Scenes via Occlusion-Aware Feature Learning}
}

%========== TECHNICAL SKILLS ==========
\section{Technical Skills}

\begin{tabularx}{\textwidth}{@{}r X@{}}
    \skillrow{Detection \& Vision}{YOLO (v5/v7/v8), Faster R-CNN, DETR, RT-DETR, Anchor-Free Detectors, Feature Pyramid Networks, NMS variants}
    \skillrow{Deep Learning}{PyTorch, CNN, Vision Transformer, Attention Mechanisms, Knowledge Distillation, Quantization-Aware Training}
    \skillrow{Deployment}{ONNX, TensorRT, OpenVINO, Triton Inference Server, CUDA, C++ inference pipelines}
    \skillrow{Engineering}{Python, C++, OpenCV, NumPy, Docker, Git, Linux, MLflow, Weights \& Biases}
    \skillrow{Data}{COCO, VOC, LVIS, Roboflow, CVAT annotation, synthetic data augmentation (Albumentations, Mosaic)}
    \skillrow{Languages}{Mandarin (native), English (fluent), German (A2)}
\end{tabularx}

%========== RESEARCH & PROJECTS ==========
\section{Research \& Engineering Projects}

\cvproject{EfficientDet-SA --- Scale-Aware Attention for Real-Time Object Detection}{2024}{%
    \begin{itemize}[leftmargin=14pt, nosep, topsep=2pt, itemsep=1pt]
        \item Designed a scale-aware attention module that dynamically re-weights multi-scale features in the FPN neck, improving small-object detection by 3.2 mAP on COCO \texttt{val2017}.
        \item Introduced a lightweight channel-spatial gating unit (0.8M extra params) compatible with YOLO and EfficientDet backbones.
        \item Achieved 42.6 mAP at 58 FPS on RTX 3090 --- Pareto-optimal on the accuracy-latency frontier vs.\ YOLOv8-M and RT-DETR-R50.
    \end{itemize}
}{PyTorch, CUDA, Albumentations, Weights \& Biases, A100 GPU}

\cvproject{HybridDet --- Transformer-CNN Hybrid Detector for Industrial Defect Inspection}{2024 -- 2025}{%
    \begin{itemize}[leftmargin=14pt, nosep, topsep=2pt, itemsep=1pt]
        \item Proposed a hybrid architecture pairing a Swin-Tiny backbone with a CNN detection head for high-resolution industrial images (4K).
        \item Deployed via TensorRT FP16 on Jetson Orin; reduced inference latency from 48ms to 11ms while maintaining 91.3\% mAP on a proprietary PCB defect dataset (12 classes, 50K images).
        \item Open-sourced training pipeline and ONNX export scripts; 200+ GitHub stars.
    \end{itemize}
}{PyTorch, TensorRT, ONNX, Jetson Orin, Docker}

\cvproject{Autonomous Driving Perception --- 3D Object Detection Module}{2023 -- 2024}{%
    \begin{itemize}[leftmargin=14pt, nosep, topsep=2pt, itemsep=1pt]
        \item Contributed to a LiDAR-camera fusion pipeline for 3D vehicle detection; implemented BEV feature alignment module.
        \item Improved nuScenes NDS by 1.8 points over BEVFusion baseline through temporal feature aggregation.
        \item Built automated evaluation CI/CD with MLflow tracking and nightly regression on 8$\times$A100 cluster.
    \end{itemize}
}{MMDetection3D, BEV Fusion, nuScenes, MLflow, Kubernetes}

\cvproject{Real-Time PPE Compliance Detection System --- Construction Sites}{2023}{%
    \begin{itemize}[leftmargin=14pt, nosep, topsep=2pt, itemsep=1pt]
        \item Built an end-to-end safety helmet and vest detection system processing 16 RTSP camera streams simultaneously.
        \item Fine-tuned YOLOv8-S on 8K annotated images; achieved 94.7\% mAP@0.5 with $<$15ms per frame on RTX 3060.
        \item Deployed as a Docker microservice with REST API; integrated with client's alert dashboard.
    \end{itemize}
}{YOLOv8, OpenCV, FastAPI, Docker, Redis}

%========== FUNDED RESEARCH ==========
\section{Funded Research Participation}

\cventry{Small Object Detection in UAV Imagery Under Complex Backgrounds}{Provincial S\&T Project}{%
    Grant No.: XXXX-2024-XXXXX\enspace$|$\enspace Role: Core Participant\par
    {\small Multi-scale feature fusion and context-aware attention for drone-captured aerial object detection.}
}

\cventry{Vision Foundation Models for Industrial Quality Inspection}{National Youth Fund}{%
    Grant No.: 62XXXXXXXX\enspace$|$\enspace Role: Participant\par
    {\small Pre-training and fine-tuning large vision models for zero-shot and few-shot defect detection.}
}

%========== PROFESSIONAL EXPERIENCE ==========
\section{Professional Experience}

\cventry{Computer Vision Engineer --- Part-Time Contract}{2024 -- Present}{%
    Developed and maintained object detection pipelines for a logistics automation startup. Responsibilities include model training, A/B testing detection architectures, and deploying TensorRT-optimized models on edge devices. Reduced package misclassification rate by 37\% through a custom cascade detection strategy.
}

%========== AWARDS ==========
\section{Awards \& Recognition}

\begin{itemize}[leftmargin=14pt, nosep, topsep=2pt, itemsep=3pt]
    \item \textbf{Gold Award} --- Provincial ``Internet+'' Innovation \& Entrepreneurship Competition, 2024.
    \item \textbf{Silver Award} ($\times$2) --- Provincial ``Internet+'' Innovation \& Entrepreneurship Competition, 2023.
    \item \textbf{Second Prize} --- National Artificial Intelligence Challenge (Object Detection Track), 2023.
    \item \textbf{First-Class Academic Scholarship} --- University of Science and Technology, 2022--2025.
    \item \textbf{National Encouragement Scholarship} --- Ministry of Education, China.
\end{itemize}

%========== BEYOND WORK ==========
\section{Beyond Work}

Passionate about the intersection of perception and robotics --- contributing to open-source detection toolboxes and writing technical blog posts on real-time inference optimization. Active participant in Kaggle competitions (Expert tier). Planning relocation to Europe for industry R\&D roles in autonomous systems or industrial AI.

\end{document}
